\begin{helper}[The odd residual set has an external third defective edge]
Fix a defective color $\delta$ and a
\noderef{ScheideDefectiveReturnData} $R$.  Write
$i=\gammaIndex(\delta)$, $j=R.j$, and
\[
 Z=V(T)\setminus(U_{\gamma_i}\cup U_{\gamma_j}).
\]
There are an exact edge identity $f_2$ and vertices $z_2,u_2$ such that
\[
 \varphi(f_2)=\delta,\quad
 \operatorname{ends}(f_2)=\{z_2,u_2\},\quad
 z_2\in Z,\quad u_2\notin V(T),\quad
 f_2\notin\{f^0_\delta,R.f_1\}.
\]
The vertex $z_2$ is distinct from $v^0_\delta$ and $R.v_1$; the edge
$f_2$ retains color $\delta$ in each of the two recolorings used in
\noderef{ScheideDefectiveDeletedTreePair}; and
$u_2\ne u^0_\delta,R.u_1$.
\end{helper}
\begin{proof}
Put $X=\noderef{TreeSequenceVertexSet}(T)$, $i=D.\gammaIndex(\delta)$,
$j=R.j$, $U_i=\{u_{i,0},u_{i,1}\}$, and
$U_j=\{u_{j,0},u_{j,1}\}$.  The fields of
\noderef{ScheideDefectiveReturnData} give $i\ne j$.  The exact vertex
positions in \noderef{ScheideStructuralConfiguration}, inherited by
\noderef{ScheideReducedConfiguration}, and vertex injectivity in
\noderef{TreeSequence} show that the four vertices in $U_i\cup U_j$ are
distinct and belong to $X$.  The same position formulas show that every
base vertex is in $X\setminus(U_i\cup U_j)$.  Thus
$|U_i\cup U_j|=4$, and
\noderef{ScheideGlobalMaximumOrder} gives
\[
 |Z|=|X|-4\quad\hbox{and}\quad |Z|\text{ odd}.
 \tag{1}
\]

By the definition of \noderef{DefectiveColors}, the color $\delta$ has
at least two exact boundary identities at $X$, so in particular it has
one.  The same maximum-order theorem says that $X$ is
\noderef{ColoredVertexSetClosed}.  If $\delta$ were missing at any
$x\in X$, that boundary identity would violate closedness.  Hence, by
the definitions of \noderef{MissingColors} and
\noderef{PresentColors}, every $z\in Z$ has an incident exact
$\delta$-edge identity.  It has only one: two distinct such identities
would be adjacent at $z$ by \noderef{AdjacentEdges} and
\noderef{Incident}, contradicting the \noderef{IsProperEdgeColoring}
field of the original \noderef{KTriple}, supplied by
\noderef{GlobalMaximumTriple} and the globally maximum field of
\noderef{ScheideStructuralConfiguration}.

Count pairs $(z,f)$ for which $z\in Z$, $f$ is incident with $z$, and
$\varphi(f)=\delta$.  The preceding existence and uniqueness make their
number $|Z|$.  An exact identity in
\noderef{InternalEdges}$(G,Z)$ contributes two pairs: its two endpoints
are distinct by \noderef{LooplessMultigraph}.  An exact identity in
\noderef{BoundaryEdges}$(G,Z)$ contributes one pair.  Exact identities
whose endpoints both lie outside $Z$ contribute none.  Properness also
makes the endpoint fibers belonging to two distinct $\delta$-identities
disjoint.  Grouping the finite pairs by their exact edge identity gives
\[
 |Z|=2|E^{\rm int}_\delta(Z)|+
   |\noderef{BoundaryColorEdges}(G,\varphi,Z,\delta)|.
 \tag{2}
\]
By (1)--(2), the exact $\delta$-colored boundary set of $Z$ has odd
cardinality.

The fields of \noderef{ScheideDefectiveSetup} give
$\varphi(f^0_\delta)=\delta$,
$\operatorname{ends}(f^0_\delta)=\{v^0_\delta,u^0_\delta\}$, and
$u^0_\delta\notin X$.  The fields of
\noderef{ScheideDefectiveReturnData} give the analogous facts for
$R.f_1,R.v_1,R.u_1$, as well as $R.v_1\ne v^0_\delta$.
Moreover, \noderef{ScheideStructuralExitsInBase} puts $v^0_\delta$ in
the base, while the return datum puts $R.v_1$ there.  Hence both tree
endpoints lie in $Z$, and $f^0_\delta,R.f_1$ belong to the boundary set
in (2).  They are distinct: equality of the identities and of their
exact endpoint pairs would either give $R.v_1=v^0_\delta$, or identify a
tree endpoint with an endpoint outside $X$.  An odd finite set containing
these two distinct identities has cardinality at least three.  Choose a
third member $f_2$ and orient its endpoints as
$\operatorname{ends}(f_2)=\{z_2,u_2\}$, with $z_2\in Z$ and
$u_2\notin Z$.  Thus $\varphi(f_2)=\delta$ and
$f_2\notin\{f^0_\delta,R.f_1\}$.

Properness shows $z_2\ne v^0_\delta,R.v_1$: otherwise $f_2$ and the
corresponding one of $f^0_\delta,R.f_1$ would be distinct adjacent exact
identities of color $\delta$.  Since $z_2\in X$, the fixed-edge clauses
of \noderef{ScheideDefectiveTerminalComponent} and
\noderef{ScheideDefectiveSecondTerminalComponent} show that both outside
swaps fix $f_2$.  In the first switched
\noderef{TwoColorIncidenceGraph}, if the edge-node of $f_2$ lay in the
relevant \noderef{TwoColorComponent}, adjacency to the vertex-node of
$z_2$ would put that node in the same component.  The exact tree-vertex
intersection stated by \noderef{ScheideDefectiveTerminalComponent}
would then give $z_2=v^0_\delta$, a contradiction.  The same argument
with \noderef{ScheideDefectiveSecondTerminalComponent} gives
$z_2=R.v_1$ for the second component, again a contradiction.  Therefore
the definition of \noderef{TwoColorComponentSwap} shows that both
component swaps fix $f_2$.  It consequently has color $\delta$ in both
final recolorings.

The same properness argument at the other endpoint gives
$u_2\ne u^0_\delta,R.u_1$.  Also $u_2\notin U_i$: if it were in $U_i$,
the endpoint $z_2\in Z\subseteq X\setminus U_i$ would make $f_2$ the
exact crossing configuration forbidden by
\noderef{ScheideDefectiveNoCross}.  Notice that
\noderef{ScheideGlobalMaximumUnused} also says that $\delta$ is unused
on the displayed edges of the original tree; however, in the two
extensions below freshness is proved directly from boundary endpoints,
not inferred from that unused-color conclusion.

It remains to exclude $u_2\in U_j$.  Suppose this membership holds, and
take the second partial tree $S_j$ from
\noderef{ScheideDefectiveDeletedTreePair}.  Its stated vertex set is
\[
 V(S_j)=(X\setminus U_j)\cup\{u^0_\delta\}.
 \tag{3}
\]
Here $z_2\in X\setminus U_j$.  On the other hand $u_2\notin V(S_j)$:
it lies in the deleted block and is not $u^0_\delta$.  Thus the exact
endpoint pair of $f_2$ proves directly, by the definition of
\noderef{BoundaryEdges}, that $f_2$ is a boundary edge of $S_j$.
Consequently \noderef{TreeSequenceBoundaryEdgeAppendExists} supplies a
fresh exact append $S_j+(f_2,u_2)$.

We verify the \noderef{TashkinovTree} obligations for this append.  All
old displayed edges retain their witnesses because $S_j$ is already a
Tashkinov tree by \noderef{ScheideDefectiveDeletedTreePair}.  For the
new edge, $f_2$ has color $\delta$ in the second final recoloring, while
$R.v_1\in X\setminus U_j\subseteq V(S_j)$ and $\delta$ is missing there
by \noderef{ScheideDefectiveSecondTerminalComponent}.  Thus
$R.v_1$ is the required earlier missing-color witness.

Let $y_0$ be the initial base vertex.  The setup field in
\noderef{ScheideDefectiveSetup} states that $\alpha _0$ is missing at
$y_0$, and the position fields of
\noderef{ScheideStructuralConfiguration} put $y_0$ in
$X\setminus U_j$.  Since $\alpha _0\ne\gamma_j$, the displayed
hypotheses of \noderef{ScheideStructuralInternalPair} give an exact
edge $g$ of color $\alpha _0$ whose endpoints are precisely the two
vertices of $U_j$.  The outside swap fixes $g$ because it is incident
with $X$.  The distinctions
$\alpha _0\ne\gamma_j,\delta$ in
\noderef{ScheideDefectiveReturnData}, together with the definition of
\noderef{TwoColorComponentSwap}, show that the component swap also
leaves its color equal to $\alpha _0$.  Stability in
\noderef{ScheideDefectiveSecondTerminalComponent}, followed by the
unaffected-color clause of
\noderef{TwoColorComponentSwapMissingBehavior}, shows that
$\alpha _0$ remains missing at $y_0$.

Let $w_2$ be the other vertex of $U_j$, so
$U_j=\{u_2,w_2\}$ and
$\operatorname{ends}(g)=\{u_2,w_2\}$.  After the first append, $u_2$ is
inside the displayed vertex set.  Equation (3), the distinct block
positions, and $u^0_\delta\notin X$ show that $w_2$ is outside it.
Thus the endpoint pair of $g$ again proves directly that $g$ is a
boundary edge, and \noderef{TreeSequenceBoundaryEdgeAppendExists}
supplies a fresh append by $(g,w_2)$.  The old Tashkinov witnesses persist,
and the new edge uses the earlier vertex $y_0$ missing $\alpha _0$;
hence this second enlarged sequence is a Tashkinov tree as well.

The coloring-independent fields of the original \noderef{KTriple},
together with the properness and exactly-one-uncolored conclusions of
\noderef{ScheideDefectiveSecondTerminalComponent}, give a $K$-triple for
the second final recoloring.  By (3), after appending both vertices of
$U_j$ the final displayed vertex set is exactly
$X\cup\{u^0_\delta\}$.  Since $u^0_\delta\notin X$, it has $|X|+1$
vertices.  Its just-verified Tashkinov property would therefore give a
candidate larger than \noderef{TashkinovOrder}$(G)$, while
\noderef{ScheideGlobalMaximumOrder} states
$|X|=\noderef{TashkinovOrder}(G)$.  This contradiction proves
$u_2\notin U_j$.

Finally, $u_2\notin Z$ by the chosen boundary orientation, and we have
proved $u_2\notin U_i\cup U_j$.  Since
$X=Z\cup U_i\cup U_j$, it follows that $u_2\notin X$.  Together with
the endpoint, color, distinctness, and recoloring facts above, this is
every assertion of the lemma.
\end{proof}
